\documentclass[11pt, oneside]{article} 
\usepackage{geometry}
\geometry{letterpaper} 
\usepackage{graphicx}
	
\usepackage{amssymb}
\usepackage{amsmath}
\usepackage{parskip}
\usepackage{color}
\usepackage{hyperref}

\graphicspath{{/Users/telliott/Dropbox/Github-Math/figures/}}
% \begin{center} \includegraphics [scale=0.4] {gauss3.png} \end{center}

\title{Euclid's Elements}
\date{}

\begin{document}
\maketitle
\Large

%[my-super-duper-separator]

\label{sec:angle_bisector}

Here's a classic problem:  can we say anything about the ratio of sides when an internal line is drawn in a triangle?  We first consider the case of a right triangle where the internal line is an angle bisector (left panel, below).  We can prove an important theorem.

\begin{center} \includegraphics [scale=0.14] {angle_bisector_b.png} \end{center}

\subsection*{angle bisector theorem}

\label{sec:angle_bisector_theorem}

\begin{center} \includegraphics [scale=0.15] {angle_bisector_c.png} \end{center}

The bisector forms two new internal triangles.  In the upper one (blue), we draw a line perpendicular to side $b$ which meets the bisector on the side opposite.  This forms two triangles with two angles the same, $\theta$ and the right angle (thus three angles the same), and one side shared.  So we have congruent right triangles by ASA.  

Thus, the two sides labeled $a$ are equal, as are the two sides labeled $c$.

The triangle with sides $c,e,d$ is similar to the whole original triangle, because both are right triangles and the top vertex is shared.  We have that the ratio of the short side to the hypotenuse is
\[ \frac{a}{b} = \frac{c}{d} \]
\[ ad = bc \]
\[ \frac{a}{c} = \frac{b}{d} \]

Tthe sides flanking the duplicated angle are in the same proportion as the parts of the base:

$\square$

\begin{center} \includegraphics [scale=0.15] {angle_bisector_c.png} \end{center}


\subsection*{bisector theorem}

The section above uses a right triangle.  However, the result is generally true.  It is called the angle bisector theorem.  Euclid's version is VI.3.

Here is a proof of the general angle bisector theorem by similar triangles.  These involve constructions that either extend the bisector to a line parallel to one side, or extend one side to meet a line parallel to the bisector.  Here we have a parallel to one side.

\begin{center} \includegraphics [scale=0.16] {angle_bisector_r7d.png} \end{center}

Given that the angle at $C$ is bisected.  The claim is that $a/y=b/x$.

\emph{Proof}.

Draw $AE$ parallel to $BC$ and extend the angle bisector to meet it at $E$.   Then the angle at $E$ is equal to the half-angle, by alternate interior angles, 

\begin{center} \includegraphics [scale=0.16] {angle_bisector_r7e.png} \end{center}

We have that $\triangle ACE$ is isosceles, so $b = b'$.

We also have two similar triangles with $\triangle AED \sim \triangle BCD$, since the angles at $D$ are equal by vertical angles.

Form the ratios of the sides opposite vertical angles to the sides opposite the angles marked with red dots:
\[ \frac{a}{y} = \frac{b'}{x} \]

But since $b = b'$:
\[ \frac{a}{y} = \frac{b}{x} \]

$\square$

\subsection*{a different view}

This looks like Euclid's proof but it's different because we don't talk about parallels, and we form a new isosceles triangle but it doesn't involve an original side.

Instead, extend the bisector of the angle between $a$ and $b$ to find a point on the circumcircle.

\begin{center} \includegraphics [scale=0.35] {bisector8.png} \end{center}

\emph{Proof}.

By equal angles on equal arcs we find the two marked angles equal to the bisected halves.  We have other equal angles, also by equal arcs.

We must be careful to choose the correct side to put in the ratios.  Recall that for similar triangles in a circle we go in different directions for the two triangles.

Start with the side opposite the vertical angle and go to the blue dotted angle, CCW for $b,m,y$ and CW for $d,x,n$.  We can pick up $b$ and $x$ with

\[ \frac{m}{x} = \frac{b}{d} \]

For the other one it's CW for $a,m,x$ and CCW for  $d,y,n$.  We pick up $a$ and $y$:
\[ \frac{m}{a} = \frac{y}{d} \]

Rearranging:
\[ md = bx = ay \]
\[ \frac{a}{b} = \frac{x}{y} \]

$\square$

As mentioned above, this theorem can be proved with a line parallel to one side and crossing the extension of the bisector, or with a line parallel to the bisector and crossing the extension of a side, or as we did in the circle above, with no parallels.

\begin{center} \includegraphics [scale=0.3] {bisector11.png} \end{center}

The figure shows both ideas with parallels.  The blue triangle extends the bisector to meet a line parallel to one of the sides originating from the other vertex.

The green triangle extends a side to meet a line parallel to the bisector.  In the second case the ratios from similar triangles are obvious.

\subsection*{external angle}

The bisector of an external angle also has a ratio.  The bisector is extended to the extension of a side.  Draw a line internal to the triangle and parallel to the bisector from $C$.  It's not labeled in the figure, but let us call it $X$.

\begin{center} \includegraphics [scale=0.35] {bisector9.png} \end{center}

\emph{Proof}.

We have two similar triangles in $\triangle AXC$ and $\triangle ABF$.  From the parts we have $BX/CF$ and from the whole we have $AB/AF$ so
\[ \frac{BX}{CF} = \frac{AB}{AF} \]

From the equal angles formed by the bisector we can find equal angles at the base of $\triangle BXC$.  It follows that the triangle is isosceles, with $BX = BC$ so then 
\[ \frac{BC}{CF} = \frac{AB}{AF} \]

$\square$

\subsection*{problem from Posamentier}

\begin{center} \includegraphics [scale=0.4] {Posamentier_mod_1.png} \end{center}

Let $O$ be the center of the circumcircle of $\triangle ABC$.

Let $CD \perp AB$ and $CX$ bisect $C$ ($\angle ACB$).

Show that $CD$ bisects $\angle ACX$.

\emph{Proof}.

Given $\angle ACX = \angle BCX$.  Draw $OC$.

The central $\angle COB$ is 2 $\angle CAB$.

But $\triangle COB$ is isosceles with equal base angles.

Thus $\angle CAB$ and $\angle BCO$ are complementary.

$\angle ACD$ is also complementary to $\angle CAB$, hence $\angle CAB$ is equal to $\angle BCO$.

\begin{center} \includegraphics [scale=0.4] {Posamentier_mod_1.png} \end{center}

Subtracting equals:  $\angle DCX = \angle OCX$.

$\square$

The angle bisector at $C$ also bisects the angle between the altitude and the radius from $C$.

\emph{Proof}.  (Alternate).

I have redrawn the figure slightly.

\begin{center} \includegraphics [scale=0.4] {Posamentier_mod_2.png} \end{center}

Claim: the altitude and diameter of the circumcircle through $C$ form equal angles with the sides $CA$ and $CB$.

Let $\triangle ABC$ have its circumcircle on center $O$, so $CO$ is a radius.  Extend $CO$ to form the diameter $CF$.

Let $CD$ be the altitude from $C$, so it cuts $AB$ at a right angle and extends to the circle at $E$.

By Thales' circle theorem, $\angle CEF$ is right.  It follows that $EF \parallel AB$.

Parallel chords in a circle cut equal arcs.  (\emph{Proof}.  Connect the opposing ends of the two chords to give equal inscribed angles.  $\square$).

Hence $AE = BF$ and then by inscribed angles, $\angle ACE = \angle BCF$.

$\square$

\begin{center} \includegraphics [scale=0.40] {Posamentier_mod_2.png} \end{center}

By subtraction, the bisector also bisects the angle between the altitude and the radius AO.

Alternatively, find $Y$ as the midpoint of arc $AB$.  

Then we have equal arcs $AY = BY$, $AE = BF$, and $EY = FY$.  Results about the angles follow easily.

\end{document}